\documentclass[10pt,a4paper]{amsart}
\usepackage{amsmath,amssymb,amsthm,graphicx}
% Retained preamble for the Stacks Project chapter snapshot in this directory.
% The chapter's own preamble.tex is a copy of the shared Stacks Project preamble; the block
% below is sliced MECHANICALLY, line for line, from it. Nothing was added, removed or
% redefined; the wrapper class, the xr-hyper/hyperref wiring and the externaldocument list
% are omitted because the assembler supplies the document class.
%
% lines taken from preamble.tex: 3, 5, 9, 11, 12, 156, 157, 158, ... (51 lines total)
%
% package coverage restored from the chapter's own preamble (2026-10-13): xy, xy-2cell

\usepackage[all]{xy}
\xyoption{2cell}
\UseAllTwocells

\newenvironment{reference}{\comment}{\endcomment}
\newenvironment{slogan}{\comment}{\endcomment}
\newenvironment{history}{\comment}{\endcomment}
\theoremstyle{plain}
\newtheorem{theorem}[subsection]{Theorem}
\newtheorem{proposition}[subsection]{Proposition}
\newtheorem{lemma}[subsection]{Lemma}
\theoremstyle{definition}
\newtheorem{definition}[subsection]{Definition}
\newtheorem{example}[subsection]{Example}
\newtheorem{exercise}[subsection]{Exercise}
\newtheorem{situation}[subsection]{Situation}
\theoremstyle{remark}
\newtheorem{remark}[subsection]{Remark}
\newtheorem{remarks}[subsection]{Remarks}
\numberwithin{equation}{subsection}
\def\lim{\mathop{\mathrm{lim}}\nolimits}
\def\colim{\mathop{\mathrm{colim}}\nolimits}
\def\Spec{\mathop{\mathrm{Spec}}}
\def\Hom{\mathop{\mathrm{Hom}}\nolimits}
\def\Ext{\mathop{\mathrm{Ext}}\nolimits}
\def\SheafHom{\mathop{\mathcal{H}\!\mathit{om}}\nolimits}
\def\SheafExt{\mathop{\mathcal{E}\!\mathit{xt}}\nolimits}
\def\Sch{\mathit{Sch}}
\def\Mor{\mathop{\mathrm{Mor}}\nolimits}
\def\Ob{\mathop{\mathrm{Ob}}\nolimits}
\def\Sh{\mathop{\mathit{Sh}}\nolimits}
\def\NL{\mathop{N\!L}\nolimits}
\def\CH{\mathop{\mathrm{CH}}\nolimits}
\def\proetale{{pro\text{-}\acute{e}tale}}
\def\etale{{\acute{e}tale}}
\def\QCoh{\mathit{QCoh}}
\def\Ker{\mathop{\mathrm{Ker}}}
\def\Im{\mathop{\mathrm{Im}}}
\def\Coker{\mathop{\mathrm{Coker}}}
\def\Coim{\mathop{\mathrm{Coim}}}
\DeclareMathSymbol{\boxtimes}{\mathbin}{AMSa}{"02}
\def\QCohstack{\mathcal{QC}\!\mathit{oh}}
\def\Cohstack{\mathcal{C}\!\mathit{oh}}
\def\Spacesstack{\mathcal{S}\!\mathit{paces}}
\def\Quotfunctor{\mathrm{Quot}}
\def\Hilbfunctor{\mathrm{Hilb}}
\def\Curvesstack{\mathcal{C}\!\mathit{urves}}
\def\Polarizedstack{\mathcal{P}\!\mathit{olarized}}
\def\Complexesstack{\mathcal{C}\!\mathit{omplexes}}
\def\Pic{\mathop{\mathrm{Pic}}\nolimits}
\def\Picardstack{\mathcal{P}\!\mathit{ic}}
\def\Picardfunctor{\mathrm{Pic}}
\def\Deformationcategory{\mathcal{D}\!\mathit{ef}}
\usepackage{cleveref}
\crefname{originallemma}{original lemma}{original lemmas}
\Crefname{originallemma}{Original lemma}{Original lemmas}
\makeatletter
\newlabel{algebraic-lemma-map-space-into-stack}{{04T4}{source}{}{sourcepointer-1}{}}
\newlabel{algebraic-lemma-map-space-into-stack@cref}{{[originallemma][0][]04T4}{[1][0][]source}{}{}{}}
\makeatother
\crefname{originallemma}{original lemma}{original lemmas}
\Crefname{originallemma}{Original lemma}{Original lemmas}
\makeatletter
\newlabel{lemma-improve-unique-point}{{lemma-improve-unique-point}{source}{}{sourcepointer-2}{}}
\newlabel{lemma-improve-unique-point@cref}{{[originallemma][0][]lemma-improve-unique-point}{[1][0][]source}{}{}{}}
\makeatother
\crefname{originallemma}{original lemma}{original lemmas}
\Crefname{originallemma}{Original lemma}{Original lemmas}
\makeatletter
\newlabel{lemma-residual-gerbe-functorial}{{lemma-residual-gerbe-functorial}{source}{}{sourcepointer-3}{}}
\newlabel{lemma-residual-gerbe-functorial@cref}{{[originallemma][0][]lemma-residual-gerbe-functorial}{[1][0][]source}{}{}{}}
\makeatother
\crefname{originallemma}{original lemma}{original lemmas}
\Crefname{originallemma}{Original lemma}{Original lemmas}
\makeatletter
\newlabel{lemma-residual-gerbe-unique}{{lemma-residual-gerbe-unique}{source}{}{sourcepointer-4}{}}
\newlabel{lemma-residual-gerbe-unique@cref}{{[originallemma][0][]lemma-residual-gerbe-unique}{[1][0][]source}{}{}{}}
\makeatother
\title{Residual gerbes map isomorphically when the field-point self-product does not change}
\author{The Stacks Project Authors}
\date{stacks-project:stacks-properties.tex:local-snapshot-20261012; GNU FDL 1.2 or later}
\begin{document}
\maketitle
\noindent\textit{Editorial note.} Statement span is lines 3125-3140 of sources/stacks-stacks-properties/pinned-original-source.tex; proof span is lines 3142-3167, the complete author proof (24 body lines).\par\medskip
\noindent\textit{Editorial note.} Navigation stubs: 4 cited-but-not-included labels are supplied as non-load-bearing pointers with their pinned tags where the snapshot tag index has them.\par\medskip
\noindent\textit{Editorial note.} extra\_packages is empty on purpose: the retained preamble already loads xy with options and a second declaration clashes in the assembler wrapper.\par\medskip
\noindent\textit{Editorial note.} Text is retained verbatim, including the author's own typos and placeholders.\par\medskip
\noindent\textit{Editorial note.} Queue index 20 of queue\_stacks-properties.json (pb=1011, stmt 3125-3140, proof 3142-3167).\par\medskip
\noindent\textit{Editorial note.} stacks-properties wave-1: fresh adjudication accept (2026-10-14); pool scanned while CONTROL.md is down (0 bytes) - §2 row deferred per Chief, claim reported to Chief 08:0x.\par\medskip
\noindent\textit{Allowed standard prerequisites (editorial).}\begin{itemize}
\item Explicit prerequisite boundary. The cited inputs are prerequisites; the argument in the retained span is not replaced by any cited result.
\end{itemize}\medskip
\ifdefined\hypertarget\hypertarget{sourcepointer-1}{}\fi
\noindent\textit{Source reference.} Original lemma 04T4: sources/stacks-stacks-properties/pinned-original-source.tex, source lines 3167--3167. Cited by the retained statement or proof but not part of the collected spans; supplied only as a navigation pointer. The argument does not depend on its content.\par\medskip
\ifdefined\hypertarget\hypertarget{sourcepointer-2}{}\fi
\noindent\textit{Source reference.} Original lemma lemma-improve-unique-point: sources/stacks-stacks-properties/pinned-original-source.tex, source lines 2791--2802. Cited by the retained statement or proof but not part of the collected spans; supplied only as a navigation pointer. The argument does not depend on its content.\par\medskip
\ifdefined\hypertarget\hypertarget{sourcepointer-3}{}\fi
\noindent\textit{Source reference.} Original lemma lemma-residual-gerbe-functorial: sources/stacks-stacks-properties/pinned-original-source.tex, source lines 3101--3113. Cited by the retained statement or proof but not part of the collected spans; supplied only as a navigation pointer. The argument does not depend on its content.\par\medskip
\ifdefined\hypertarget\hypertarget{sourcepointer-4}{}\fi
\noindent\textit{Source reference.} Original lemma lemma-residual-gerbe-unique: sources/stacks-stacks-properties/pinned-original-source.tex, source lines 3076--3086. Cited by the retained statement or proof but not part of the collected spans; supplied only as a navigation pointer. The argument does not depend on its content.\par\medskip
\begin{lemma}
\label{lemma-residual-gerbe-isomorphic}
Let $f : \mathcal{X} \to \mathcal{Y}$ be a morphism of algebraic stacks.
Let $x \in |\mathcal{X}|$ with image $y \in |\mathcal{Y}|$.
Assume the residual gerbes $\mathcal{Z}_x \subset \mathcal{X}$
and $\mathcal{Z}_y \subset \mathcal{Y}$ of $x$ and $y$ exist
and that there exists a morphism $\Spec(k) \to \mathcal{X}$
in the equivalence class of $x$ such that
$$
\Spec(k) \times_\mathcal{X} \Spec(k)
\longrightarrow
\Spec(k) \times_\mathcal{Y} \Spec(k)
$$
is an isomorphism. Then $\mathcal{Z}_x \to \mathcal{Z}_y$
is an isomorphism.
\end{lemma}

\begin{proof}
Let $k'/k$ be an extension of fields. Then
$$
\Spec(k') \times_\mathcal{X} \Spec(k')
\longrightarrow
\Spec(k') \times_\mathcal{Y} \Spec(k')
$$
is the base change of the morphism in the lemma by the
faithfully flat morphism $\Spec(k' \otimes k') \to \Spec(k \otimes k)$.
Thus the property described in the lemma is independent of the
choice of the morphism $\Spec(k) \to \mathcal{X}$ in the
equivalence class of $x$. Thus we may assume that
$\Spec(k) \to \mathcal{Z}_x$ is surjective, flat, and
locally of finite presentation. In this situation we have
$$
\mathcal{Z}_x = [\Spec(k)/R]
$$
with $R = \Spec(k) \times_\mathcal{X} \Spec(k)$. See
proof of Lemma \ref{lemma-improve-unique-point}.
Since also $R = \Spec(k) \times_\mathcal{Y} \Spec(k)$
we conclude that the morphism $\mathcal{Z}_x \to \mathcal{Z}_y$
of Lemma \ref{lemma-residual-gerbe-functorial}
is fully faithful by
Algebraic Stacks, Lemma \ref{algebraic-lemma-map-space-into-stack}.
We conclude for example by Lemma \ref{lemma-residual-gerbe-unique}.
\end{proof}

\end{document}
